\documentclass[11pt, letterpaper]{article}

\usepackage[utf8]{inputenc}
\usepackage[T1]{fontenc}
\usepackage{lmodern}

\usepackage{amsmath, amssymb, amsthm, amsfonts, mathrsfs, mathtools}

\usepackage{geometry}
\geometry{letterpaper, margin=1in}
\setlength{\parindent}{0pt}
\setlength{\parskip}{0.6em}

\usepackage{graphicx, xcolor, enumitem, booktabs, array, tabularx}

\begin{document}

\begin{center}
\textbf{\Large Homework 8 Chapter 5} \\
Tashfeen Hatim Said Omran \\
UCF ID: 5862894 \\
Class: STA4321 \\
Novermber 11, 2026
\end{center}


\subsection{Exercise 5.22}
\begin{enumerate}[label=\textbf{\alph*}]
    \item Give the marginal probability functions for $Y_1$ and $Y_2$.
    \item Give the conditional probability function for $Y_2$ given $Y_1 = 0$.
    \item What is the probability that a child survived given that he or she was in a car-seat belt?
\end{enumerate}

\textbf{Solution:}

\textbf{Part (a):} By Definition 5.4(a), reading the column and row totals directly from the joint table gives:
\begin{itemize}
    \item \textbf{Marginal PMF of $Y_1$:} $p_1(0) = 0.76$, $p_1(1) = 0.24$, and $p_1(y_1) = 0$ elsewhere.
    \item \textbf{Marginal PMF of $Y_2$:} $p_2(0) = 0.55$, $p_2(1) = 0.16$, $p_2(2) = 0.29$, and $p_2(y_2) = 0$ elsewhere.
\end{itemize}

\textbf{Part (b):} By Definition 5.5, for each $y_2 \in \{0, 1, 2\}$, $p(y_2 \mid Y_1 = 0) = \frac{p(0, y_2)}{p_1(0)} = \frac{p(0, y_2)}{0.76}$:
\begin{align*}
    p(Y_2 = 0 \mid Y_1 = 0) &= \frac{0.38}{0.76} = \frac{1}{2} = 0.50, \\
    p(Y_2 = 1 \mid Y_1 = 0) &= \frac{0.14}{0.76} = \frac{7}{38} \approx 0.1842, \\
    p(Y_2 = 2 \mid Y_1 = 0) &= \frac{0.24}{0.76} = \frac{6}{19} \approx 0.3158,
\end{align*}
and $p(y_2 \mid Y_1 = 0) = 0$ elsewhere.

\textbf{Part (c):} Here what i want is$\Prob(\text{child survived} \mid \text{car-seat belt used}) = \Prob(Y_1 = 0 \mid Y_2 = 2)$. By Definition 5.5:
\[ \Prob(Y_1 = 0 \mid Y_2 = 2) = \frac{p(0, 2)}{p_2(2)} = \frac{0.24}{0.29} = \frac{24}{29} \approx 0.8276. \quad \blacksquare \]

\subsection{Exercise 5.25}

\textbf{Solution:}

\textbf{Part (a):} For $y_1 > 0$:
\[ f_1(y_1) = \int_0^\infty e^{-y_1} e^{-y_2} \, dy_2 = e^{-y_1} \left[ -e^{-y_2} \right]_0^\infty = e^{-y_1}, \quad y_1 > 0, \]
and $f_1(y_1) = 0$ elsewhere. By symmetry, $f_2(y_2) = e^{-y_2}$ for $y_2 > 0$, and $0$ elsewhere. Both $Y_1$ and $Y_2$ follow an \textbf{Exponential distribution} with mean $\beta = 1$ (or a Gamma distribution :) with $\alpha = 1, \beta = 1$).

\textbf{Part (b):} Integrating the marginal densities from $1$ to $2.5$:
\[ \Prob(1 < Y_1 < 2.5) = \int_1^{2.5} e^{-y_1} \, dy_1 = e^{-1} - e^{-2.5} \approx 0.2858, \qquad \Prob(1 < Y_2 < 2.5) = e^{-1} - e^{-2.5} \approx 0.2858. \]

\textbf{Part (c):} Since $f_2(y_2) = e^{-y_2} > 0$ for all $y_2 > 0$ and $f_2(y_2) = 0$ for $y_2 \le 0$, $f(y_1 \mid y_2)$ is defined for $y_2 > 0$.

\textbf{Part (d):} For any $y_2 > 0$:
\[ f(y_1 \mid y_2) = \frac{f(y_1, y_2)}{f_2(y_2)} = \begin{cases} \dfrac{e^{-y_1}e^{-y_2}}{e^{-y_2}} = e^{-y_1}, & y_1 > 0, \\ 0, & \text{elsewhere.} \end{cases} \]

\textbf{Part (e):} For any $y_1 > 0$:
\[ f(y_2 \mid y_1) = \frac{f(y_1, y_2)}{f_1(y_1)} = \begin{cases} \dfrac{e^{-y_1}e^{-y_2}}{e^{-y_1}} = e^{-y_2}, & y_2 > 0, \\ 0, & \text{elsewhere.} \end{cases} \]

\textbf{Part (f):} For every $y_2 > 0$, the conditional density $f(y_1 \mid y_2)$ is \textbf{identical} to the marginal density $f_1(y_1)$ found in part (a): $f(y_1 \mid y_2) = f_1(y_1) = e^{-y_1} \mathbf{1}_{(0, \infty)}(y_1)$.

\textbf{Part (g):} Because $f(y_1 \mid y_2) = f_1(y_1)$ for all $y_1 \in \mathbb{R}$ and all $y_2 > 0$, integrating over any interval $I = (a, b) \subset \mathbb{R}$ yields:
\[ \Prob(a < Y_1 < b \mid Y_2 = y_2) = \int_a^b f(y_1 \mid y_2) \, dy_1 = \int_a^b f_1(y_1) \, dy_1 = \Prob(a < Y_1 < b). \]
Thus, the marginal and conditional probabilities that $Y_1$ falls in any interval are equal, proving that $Y_1$ and $Y_2$ are statistically independent. $\blacksquare$

\subsection{Exercise 5.33}

\textbf{Solution:}

\textbf{Part (a):} For $0 \le y_1 < \infty$, integrating with respect to $y_2 \in [0, y_1]$:
\[ f_1(y_1) = \int_0^{y_1} e^{-y_1} \, dy_2 = y_1 e^{-y_1}, \quad y_1 \ge 0, \text{ and } 0 \text{ elsewhere.} \]
For $0 \le y_2 < \infty$, integrating with respect to $y_1 \in [y_2, \infty)$:
\[ f_2(y_2) = \int_{y_2}^\infty e^{-y_1} \, dy_1 = \left[ -e^{-y_1} \right]_{y_2}^\infty = e^{-y_2}, \quad y_2 \ge 0, \text{ and } 0 \text{ elsewhere.} \]

\textbf{Part (b):} Since $f_2(y_2) = e^{-y_2} > 0$ for all $y_2 \ge 0$, $f(y_1 \mid y_2)$ is defined for all $y_2 \ge 0$:
\[ f(y_1 \mid y_2) = \frac{f(y_1, y_2)}{f_2(y_2)} = \begin{cases} \dfrac{e^{-y_1}}{e^{-y_2}} = e^{-(y_1 - y_2)}, & y_1 \ge y_2, \\ 0, & \text{elsewhere.} \end{cases} \]

\textbf{Part (c):} Since $f_1(y_1) = y_1 e^{-y_1} > 0$ if and only if $y_1 > 0$, $f(y_2 \mid y_1)$ is defined for all $y_1 > 0$:
\[ f(y_2 \mid y_1) = \frac{f(y_1, y_2)}{f_1(y_1)} = \begin{cases} \dfrac{e^{-y_1}}{y_1 e^{-y_1}} = \dfrac{1}{y_1}, & 0 \le y_2 \le y_1, \\ 0, & \text{elsewhere.} \end{cases} \]

\textbf{Part (d):} \textbf{No.} The conditional density $f(y_1 \mid y_2) = e^{-(y_1 - y_2)} \mathbf{1}_{[y_2, \infty)}(y_1)$ depends explicitly on $y_2$ (both in its formula and in its support) and is not equal to the marginal density $f_1(y_1) = y_1 e^{-y_1} \mathbf{1}_{[0, \infty)}(y_1)$.

\textbf{Part (e):} It implies that the marginal probability $\Prob(a < Y_1 < b)$ and the conditional probability $\Prob(a < Y_1 < b \mid Y_2 = y_2)$ are \textbf{not equal in general}; knowing the waiting time $Y_2 = y_2$ alters the probability distribution of the total time $Y_1$, meaning $Y_1$ and $Y_2$ are \textbf{dependent} random variables. $\blacksquare$

\subsection{Exercise 5.48}

\textbf{Solution:} 

Are $Y_1$ (child survived) and $Y_2$ (belt used) independent?
\begin{proof}
From the provided table, extract the joint and marginal probabilities for the cell $(0,0)$:
\[ p(0,0) = 0.38 \]
\[ p_1(0) = 0.55 \quad \text{and} \quad p_2(0) = 0.76 \]
I'll then check the factorization condition:
\[ p_1(0)p_2(0) = 0.55 \times 0.76 = 0.418 \]
Since $0.38 \neq 0.418$, the variables are \textbf{dependent}. $\blacksquare$
\end{proof}

\subsection{Exercise 5.51}

\textbf{Solution:} 

$f(y_1, y_2) = e^{-(y_1+y_2)}$ for $y_1 > 0, y_2 > 0$.
\begin{proof}
The support is $(0, \infty) \times (0, \infty)$, which has constant boundaries. The density factors as:
\[ f(y_1, y_2) = e^{-y_1} e^{-y_2} \]
Let $g(y_1) = e^{-y_1}$ and $h(y_2) = e^{-y_2}$. By Theorem 5.5, $Y_1$ and $Y_2$ are \textbf{independent}. $\blacksquare$
\end{proof}

\subsection{Problem 5.76}
\textbf{Solution:}
The variables are independent since $f(y_1, y_2) = (2y_1)(2y_2)$. The marginal is $f_1(y_1) = 2y_1$ for $y_1 \in [0,1]$.
\begin{enumerate}[label=\alph*)]
    \item $\E[Y_1] = \int_0^1 y_1 (2y_1) \, dy_1 = \int_0^1 2y_1^2 \, dy_1 = \left[ \frac{2}{3}y_1^3 \right]_0^1 = \frac{2}{3}$.
    \item $\E[Y_1^2] = \int_0^1 y_1^2 (2y_1) \, dy_1 = \int_0^1 2y_1^3 \, dy_1 = \left[ \frac{1}{2}y_1^4 \right]_0^1 = \frac{1}{2}$. \\ \relax
    $\Var(Y_1) = \E[Y_1^2] - (\E[Y_1])^2 = \frac{1}{2} - \left(\frac{2}{3}\right)^2 = \frac{1}{2} - \frac{4}{9} = \frac{9 - 8}{18} = \frac{1}{18}$.
    \item By symmetry, $\E[Y_2] = 2/3$. Thus, $\E[Y_1 - Y_2] = \frac{2}{3} - \frac{2}{3} = 0$.
\end{enumerate}
\hfill $\blacksquare$


\subsection{Problem 5.77}
\textbf{question:} $f(y_1, y_2) = 6(1 - y_2)$ for $0 \leq y_1 \leq y_2 \leq 1$. Find expectations and variances.

\textbf{solution:}
finding the the marginal densities first.
For $Y_2$:
\[
f_2(y_2) = \int_0^{y_2} 6(1 - y_2) \, dy_1 = 6(1 - y_2)[y_1]_0^{y_2} = 6y_2(1 - y_2), \quad 0 \leq y_2 \leq 1
\]
This is a Beta distribution with $\alpha=2, \beta=2$.
For $Y_1$:
\[
f_1(y_1) = \int_{y_1}^1 6(1 - y_2) \, dy_2 = 6 \left[ y_2 - \frac{y_2^2}{2} \right]_{y_1}^1 = 6 \left( \frac{1}{2} - y_1 + \frac{y_1^2}{2} \right) = 3(1 - y_1)^2, \quad 0 \leq y_1 \leq 1
\]
and this is a Beta distribution with $\alpha=1, \beta=3$.
\begin{enumerate}[label=\alph*)]
    \item $\E[Y_1] = \int_0^1 3y_1(1-y_1)^2 \, dy_1 = 3 \int_0^1 (y_1 - 2y_1^2 + y_1^3) \, dy_1 = 3\left(\frac{1}{2} - \frac{2}{3} + \frac{1}{4}\right) = 3\left(\frac{1}{12}\right) = \frac{1}{4}$. \\ \relax
    $\E[Y_2] = \int_0^1 6y_2^2(1-y_2) \, dy_2 = 6\left(\frac{1}{3} - \frac{1}{4}\right) = 6\left(\frac{1}{12}\right) = \frac{1}{2}$.
    \item $\E[Y_1^2] = \int_0^1 3y_1^2(1-y_1)^2 \, dy_1 = 3\left(\frac{1}{3} - \frac{2}{4} + \frac{1}{5}\right) = 3\left(\frac{1}{30}\right) = \frac{1}{10}$. \\ \relax
    $\Var(Y_1) = \frac{1}{10} - \left(\frac{1}{4}\right)^2 = \frac{1}{10} - \frac{1}{16} = \frac{8 - 5}{80} = \frac{3}{80}$. \\ \relax
    $\E[Y_2^2] = \int_0^1 6y_2^3(1-y_2) \, dy_2 = 6\left(\frac{1}{4} - \frac{1}{5}\right) = \frac{6}{20} = \frac{3}{10}$. \\ \relax
    $\Var(Y_2) = \frac{3}{10} - \left(\frac{1}{2}\right)^2 = \frac{3}{10} - \frac{1}{4} = \frac{6 - 5}{20} = \frac{1}{20}$.
    \item $\E[Y_1 - 3Y_2] = \E[Y_1] - 3\E[Y_2] = \frac{1}{4} - 3\left(\frac{1}{2}\right) = \frac{1}{4} - \frac{6}{4} = -\frac{5}{4}$.
\end{enumerate}
\hfill $\blacksquare$



\end{document}